La hipótesis o definición del problema --según la modalidad de trabajo de graduación-- es una respuesta tentativa a la pregunta de investigación. Es una declaración que puede validarse estadísticamente o mediante información empírica. Guía la investigación puesto que establece los límites, enfoca el problema y organiza el pensamiento. Algunos ejemplos: 

\begin{itemize}
    \item Las nanopartículas de carboximetilquitosano mantendrán su carga máxima de adsorción de plomo tras tres ciclos de adsorción/desorción.

    \item Las producciones de entretenimiento —-libros, películas y series-—, distribuidas por la industria cultural, impactan los significados que los guatemaltecos entre 20-39 y 40-59 años tienen sobre la inteligencia artificial.
    
\end{itemize}

\noindent
Tome en cuenta los siguientes consejos para redactarla:
\begin{itemize}
    \item Es una oración enunciativa (sujeto + verbo + objeto) en forma impersonal (tercera persona).

    \item Debe ser una afirmación sobre algo nuevo.

    \item Indica la relación entre variables.

    \item Es comprobable y cuantificable.
\end{itemize}

Si la definición del problema realmente consiste en una delimitación (recursos, enfoque, tiempo, resultados esperados, etc.), entonces el capítulo recibe el nombre de alcance y debe estar narrativamente amarrado a las recomendaciones del trabajo.